If \(h\in V_m\), write \(h=\sum_c\gamma_cB_c^{(m)}\).  For every feasible \(b\) and every \(\nu\in\mathcal F_{m,\epsilon}(b)\),
\[
 \int h\,d\nu=\sum_c\gamma_c\int B_c^{(m)}\,d\nu
 =\sum_c\gamma_cb_c,
\]
so the functional is constant on every fiber.

Conversely, suppose \(h\notin V_m\).  Choose \(b_0\in\operatorname{relint}(\mathcal B_{m,\epsilon})\), which is nonempty for every nonempty finite-dimensional convex set.  The full-support perturbation lemma supplies \(\nu_0\in\mathcal F_{m,\epsilon}(b_0)\) with full support.  Project \(h\) onto \(V_m\) in \(L^2(\nu_0)\) and put \(r=h-\Pi_{V_m}^{L^2(\nu_0)}h\).  The same orthogonality calculation gives
\[
 \int B_c^{(m)}r\,d\nu_0=0,
 \qquad
 \int hr\,d\nu_0=\int r^2\,d\nu_0>0.
\]
The strict inequality follows because a continuous function that vanishes \(\nu_0\)-almost surely vanishes everywhere under full support.  For sufficiently small \(s>0\), \((1\pm sr)\nu_0\) are probability measures with the same moment vector \(b_0\) and different integrals of \(h\).  Thus constancy on every feasible fiber implies \(h\in V_m\), proving the equivalence.

For \(h=f\), the finite-plate nonpolynomial lemma gives \(f\notin V_m\), so universal point identification fails.  For every relative-interior \(b\), the two measures constructed by the full-support perturbation lemma lie in \(\mathcal F_{m,\epsilon}(b)\) and have distinct \(f\)-integrals.  Since \([L_m(b),U_m(b)]\) contains both, \(L_m(b)<U_m(b)\).